% Part: computability
% Chapter: computability-theory
% Section: russells-paradox

\documentclass[../../../include/open-logic-section]{subfiles}

\begin{document}

\olfileid[hi]{cmp}{thy}{rus}
\olsection{रसेल के विरोधाभास से तुलना}

इस अनुभाग के तर्कों की रसेल के विरोधाभास से समानता और भिन्नता देखना
ज्ञानवर्धक है:
\begin{enumerate}
\item रसेल का विरोधाभास: $S = \Setabs{x}{x \notin x}$ मानें। तब $S
  \in S$ तभी और केवल तभी जब $X \notin S$—यह विरोधाभास है।

  \emph{निष्कर्ष:} ऐसा कोई समुच्चय~$S$ नहीं है। ``सभी समुच्चयों के
  समुच्चय'' का अस्तित्व मानना समुच्चय सिद्धांत के अन्य अभिगृहीतों से
  असंगत है।

\item रसेल के विरोधाभास का एक रूपांतरण: $F$ को सभी फलनों के समुच्चय से
  $\{ 0, 1 \}$ में वह ``फलन'' मानें जिसकी परिभाषा है
  \[
  F(f) =
  \begin{cases}
    1 & \text{यदि $f$, $f$ के प्रांत में है और $f(f) = 0$} \\
    0 & \text{अन्यथा}
  \end{cases}
  \]
  ऐसा ही तर्क दिखाता है कि $F(F) = 0$ तभी और केवल तभी जब $F(F) = 1$—
  यह विरोधाभास है।

  \emph{निष्कर्ष:} $F$ कोई फलन नहीं है। ``सभी फलनों का समुच्चय'' इतना
  बड़ा है कि वह किसी फलन का प्रांत नहीं हो सकता।

\item विकर्णीकरण तर्क: $f_0$, $f_1$, \dots को आंशिक संगणनीय फलनों का
  परिगणन मानें, और $G \colon \Nat \to
  \{ 0, 1 \}$ की परिभाषा हो
  \[
  G(x) =
  \begin{cases}
    1 & \text{यदि $f_x(x)\downarrow = 0$} \\
    0 & \text{अन्यथा}
  \end{cases}
  \]
  यदि $G$ संगणनीय है, तो वह फलन $f_k$ है, किसी $k$ के लिए। किंतु तब
  $G(k) = 1$ तभी और केवल तभी जब $G(k) = 0$—यह विरोधाभास है।

  \emph{निष्कर्ष:} $G$ संगणनीय नहीं है। ध्यान दें कि समुच्चय सिद्धांत के
  अभिगृहीतों के अनुसार $G$ फिर भी एक फलन है; यहाँ कोई विरोधाभास नहीं,
  केवल एक स्पष्टीकरण है।
\end{enumerate}

आंशिक फलनों, संगणनीय फलनों, आंशिक संगणनीय फलनों आदि की यह चर्चा भ्रमित
कर सकती है। $\Nat$ से $\Nat$ में सभी आंशिक फलनों का समुच्चय वस्तुओं का
एक बड़ा संग्रह है। उनमें कुछ पूर्ण हैं, कुछ संगणनीय हैं, कुछ पूर्ण और
संगणनीय दोनों हैं, और कुछ इनमें से कोई नहीं हैं। ध्यान रखें कि जब हम
``फलन'' कहते हैं, तो सामान्यतः हमारा अर्थ पूर्ण फलन होता है। अतः हमारे
पास ये वर्ग हैं:
\begin{enumerate}
\item संगणनीय फलन
\item वे आंशिक संगणनीय फलन जो पूर्ण नहीं हैं
\item वे फलन जो संगणनीय नहीं हैं
\item वे आंशिक फलन जो न पूर्ण हैं, न संगणनीय
\end{enumerate}
इसे व्यवस्थित करने के लिए $\Nat$ से $\Nat$ में सभी आंशिक फलनों को
दर्शाने वाला एक बड़ा वर्ग बनाना और फिर क्रमशः पूर्ण फलनों तथा संगणनीय
आंशिक फलनों के अनुरूप दो अतिव्यापी क्षेत्र चिह्नित करना सहायक हो सकता है।
आरेख में बने प्रत्येक क्षेत्र की किसी वस्तु का वर्णन कर पाना एक अच्छा
अभ्यास है।

\end{document}
